\documentclass[11pt]{article}

\usepackage[T1]{fontenc}
\usepackage[margin=2.2cm]{geometry}
\usepackage{amsmath}
\usepackage{graphicx}
\usepackage{microtype}
\usepackage{xcolor}
\usepackage{hyperref}

\hypersetup{
  colorlinks=true,
  linkcolor=black,
  urlcolor=blue!60!black
}

\setlength{\parindent}{0pt}
\setlength{\parskip}{0.65em}
\renewcommand{\arraystretch}{1.15}

\definecolor{promptbg}{RGB}{245,246,248}
\definecolor{rulegray}{RGB}{210,213,218}

\newcommand{\assetroot}{../data/objaverse_eval/0bf9c0d7fd69472e8fb192e0a784b165}
\newcommand{\promptprefix}{\texttt{refcontrol}, preserve the reference object's
  identity, geometry, surface materials, textures, markings, and colors. Change
  only the illumination.}

\newcommand{\promptbox}[1]{%
  \begingroup
  \setlength{\fboxsep}{8pt}%
  \noindent\colorbox{promptbg}{%
    \parbox{\dimexpr\linewidth-2\fboxsep\relax}{%
      \small\textbf{Prompt}\quad \promptprefix\ #1}}
  \endgroup
}

\newcommand{\promptheader}[1]{%
  \begin{minipage}[b]{0.255\textwidth}
    \scriptsize\raggedright
    \promptprefix\ #1\par
  \end{minipage}%
}

\begin{document}

\begin{center}
  {\LARGE\bfseries Progress Update for Week 6.07--12.07}
\end{center}

\section{Multi-Illumination Images}

We generate images of the same object under different lighting conditions with
the \href{https://huggingface.co/thedeoxen/refcontrol-FLUX.2-klein-9B-reference-normal-lora}{RefControl FLUX.2 Klein 9B Reference--Normal LoRA}.
For a fixed camera view, the model receives both a rendered normal map and a
high-resolution RGB reference. The normal map constrains the geometry and
surface orientation, while the RGB image provides the object's materials,
textures, markings, and colours. A text prompt then specifies the desired
illumination. Only the lighting clause is changed between generations.

\begin{figure}[htbp]
  \centering
  \begin{tabular}{@{}c@{\hspace{0.7em}}c@{\hspace{0.7em}}c@{\hspace{0.7em}}c@{\hspace{0.7em}}c@{}}
    \includegraphics[width=0.275\textwidth]{\assetroot/renders/view_002/normals.png} &
    {\Large$+$} &
    \includegraphics[width=0.275\textwidth]{\assetroot/renders/view_002/rgb.png} &
    {\Large$\longrightarrow$} &
    \includegraphics[width=0.275\textwidth]{\assetroot/relightings/view_002/point_left.png}
    \\
    \textbf{Normal map} && \textbf{RGB reference} && \textbf{Relit image}
  \end{tabular}
  \vspace{0.7em}

  \promptbox{Single point light from the left, positioned outside the camera
  frame; no visible light source or light beam}

  \caption{Relighting workflow for a single camera view. The normal map and RGB
  reference are taken from the same view; the shown prompt produces the
  corresponding left-lit output.}
  \label{fig:relighting-workflow}
\end{figure}

Using a high-resolution RGB reference is important for appearance consistency.
When fine texture details are not sufficiently resolved, the underlying
generative model can reinterpret them independently in each output. In this
example, this can include changing the colour of the small blue accent dots on
the object. We therefore generate from $1024\times1024$ RGB and normal-map
inputs, always paired from the same camera view.

\subsection{Examples}

Figure~\ref{fig:relighting-examples} compares three illumination prompts at
camera yaw angles of $-30^\circ$, $0^\circ$, and $+30^\circ$ (all rendered at a
fixed $20^\circ$ elevation). Each column uses the prompt printed above it, while
each row shows a different view of the object. This layout makes it possible to
inspect whether a requested illumination remains consistent as the camera moves
around the object.

\begin{figure}[p]
  \centering
  \begin{tabular}{@{}r@{\hspace{0.7em}}c@{\hspace{0.55em}}c@{\hspace{0.55em}}c@{}}
    &
    \promptheader
      {Single point light from the left, positioned outside the camera frame;
       no visible light source or light beam}
    &
    \promptheader
      {Single point light from above, positioned outside the camera frame; no
       visible light source or light beam}
    &
    \promptheader
      {Large soft light from the right, positioned outside the camera frame; no
       visible light source or light beam}
    \\[0.7em]
    \textbf{View 0} ($-30^\circ$) &
    \includegraphics[width=0.255\textwidth]{\assetroot/relightings/view_000/point_left.png} &
    \includegraphics[width=0.255\textwidth]{\assetroot/relightings/view_000/point_top.png} &
    \includegraphics[width=0.255\textwidth]{\assetroot/relightings/view_000/soft_right.png}
    \\[0.55em]
    \textbf{View 2} ($0^\circ$) &
    \includegraphics[width=0.255\textwidth]{\assetroot/relightings/view_002/point_left.png} &
    \includegraphics[width=0.255\textwidth]{\assetroot/relightings/view_002/point_top.png} &
    \includegraphics[width=0.255\textwidth]{\assetroot/relightings/view_002/soft_right.png}
    \\[0.55em]
    \textbf{View 4} ($+30^\circ$) &
    \includegraphics[width=0.255\textwidth]{\assetroot/relightings/view_004/point_left.png} &
    \includegraphics[width=0.255\textwidth]{\assetroot/relightings/view_004/point_top.png} &
    \includegraphics[width=0.255\textwidth]{\assetroot/relightings/view_004/soft_right.png}
  \end{tabular}
  \caption{Three illumination conditions generated independently at camera yaw
  angles of $-30^\circ$, $0^\circ$, and $+30^\circ$. Rows correspond to camera
  views and columns to prompts.}
  \label{fig:relighting-examples}
\end{figure}

\clearpage
\section{Albedo Estimation}

An observed image mixes properties of the surface with the illumination under
which it was captured. Following the intrinsic-image terminology introduced by
\href{https://people.csail.mit.edu/torralba/courses/6.870/papers/Barrow-Tenenbaum78.pdf}{Barrow and Tenenbaum (1978)}, we model an image as the element-wise
product
\begin{equation}
  I_t(x,y) = L_t(x,y) \odot R(x,y),
  \label{eq:intrinsic-image}
\end{equation}
where $L_t$ is the illumination or shading at time $t$ and $R$ is the
reflectance (albedo). Albedo describes the colour attributable to the surface
itself and should remain invariant when only the illumination changes.

In \emph{2D albedo estimation}, the camera and object are fixed and only the
illumination varies. The task is to recover the single reflectance image shared
by the sequence, matching the problem posed by
\href{https://people.eecs.berkeley.edu/~efros/courses/AP06/Papers/weiss-iccv-01.pdf}{Weiss (2001)}. In \emph{3D albedo estimation}, observations may come from both
different illuminations and different camera views. The goal is then one
surface-aligned reflectance texture that is consistent wherever the same 3D
point is visible, as in
\href{https://people.csail.mit.edu/sparis/publi/2012/sigasia_intrinsic/Laffont_12_Coherent_Intrinsic_Images.pdf}{\emph{Coherent Intrinsic Images from Photo Collections}}.

\subsection{Naive Marginalisation}

Since reflectance is constant across the relit sequence, a simple baseline is
to marginalise over illumination by aggregating the observations. For $T$
registered relightings of one camera view, the two estimates used here are
\begin{equation}
  \widehat{R}_{\mathrm{mean}}(x,y)
    = \frac{1}{T}\sum_{t=1}^{T} I_t(x,y),
  \qquad
  \widehat{R}_{\mathrm{median}}(x,y)
    = \underset{t}{\operatorname{median}}\; I_t(x,y).
  \label{eq:naive-aggregation}
\end{equation}
The operations are applied per pixel and colour channel in linear RGB, after
masking the background and rejecting saturated observations.

\begin{figure}[htbp]
  \centering
  \begin{tabular}{@{}c@{\hspace{0.7em}}c@{\hspace{0.7em}}c@{}}
    \includegraphics[width=0.30\textwidth]{\assetroot/albedo_2d/naive_mean/view_002/albedo.png} &
    \includegraphics[width=0.30\textwidth]{\assetroot/albedo_2d/naive_median/view_002/albedo.png} &
    \includegraphics[width=0.30\textwidth]{\assetroot/renders/view_002/gt_albedo.png}
    \\
    \textbf{Mean} & \textbf{Median} & \textbf{Ground truth}
  \end{tabular}
  \caption{Naive 2D albedo estimates from 15 relightings of view 2, compared
  with the ground-truth reflectance.}
  \label{fig:naive-albedo-2d}
\end{figure}

The mean and median retain a considerable amount of shading, particularly on
the barrel and grip, instead of recovering the flatter ground-truth colours.
With unknown, non-uniform lighting and only 15 samples, the illumination does
not average to a neutral field; consequently, its bias is baked into both
albedo estimates.

\subsubsection{Aggregation in 3D}

We also test whether pooling observations across camera views reduces this
bias. Each render stores, for every visible object pixel, its coordinate in the
mesh's UV atlas. We convert the relit images to linear RGB, discard invalid and
saturated pixels, and splat the remaining colours into a shared $1024\times
1024$ UV texture. A per-texel mean or median is then taken over all valid
lighting and view observations. For comparison below, the resulting texture is
reprojected into view 2 using the same UV-coordinate buffer.

\begin{figure}[htbp]
  \centering
  \begin{tabular}{@{}c@{\hspace{0.7em}}c@{\hspace{0.7em}}c@{}}
    \includegraphics[width=0.285\textwidth]{\assetroot/albedo_3d/naive_mean/visualisations/view_002/albedo.png} &
    \includegraphics[width=0.285\textwidth]{\assetroot/albedo_3d/naive_median/visualisations/view_002/albedo.png} &
    \includegraphics[width=0.285\textwidth]{\assetroot/renders/view_002/gt_albedo.png}
    \\
    \textbf{Mean in UV space} & \textbf{Median in UV space} & \textbf{Ground truth}
  \end{tabular}
  \caption{Naive 3D albedo estimates aggregated in a common UV atlas and
  reprojected into view 2.}
  \label{fig:naive-albedo-3d}
\end{figure}

Aggregating in UV space helps to suppress some view-dependent specular
reflections. Nevertheless, both estimates still contain substantial shading,
so adding observations from multiple views does not resolve the underlying
illumination bias.

\subsection{Gradient Space Marginalisation}

The method of
\href{https://people.eecs.berkeley.edu/~efros/courses/AP06/Papers/weiss-iccv-01.pdf}{Weiss, \emph{Deriving Intrinsic Images from Image Sequences} (2001)}
performs the marginalisation in gradient space. Taking logarithms turns the
image formation model into a sum,
\begin{equation}
  \log I_t = \log L_t + \log R.
  \label{eq:log-intrinsic-image}
\end{equation}
Spatial gradients caused by illumination are assumed to be sparse and to move
between differently lit images, whereas gradients at reflectance boundaries
remain in the same location. We apply horizontal and vertical difference
filters to each log image and take the temporal median of their responses.
Under these assumptions, the moving illumination gradients are rejected while
the persistent reflectance gradients remain. A masked Poisson solve then
integrates the recovered gradient field into an albedo image.

The same construction can be applied in 3D. We first project every relighting
from every view into the shared UV atlas, then compute and median-aggregate the
log-RGB gradients in UV space. Only gradients supported by multiple
observations and camera views are retained, and a masked Poisson solve recovers
one surface-aligned albedo texture. The texture is reprojected into view 2 for
the comparison below.

\clearpage
\begin{figure}[htbp]
  \centering
  \begin{tabular}{@{}c@{\hspace{0.7em}}c@{\hspace{0.7em}}c@{}}
    \includegraphics[width=0.30\textwidth]{\assetroot/albedo_2d/weiss/view_002/albedo.png} &
    \includegraphics[width=0.30\textwidth]{\assetroot/albedo_3d/weiss/visualisations/view_002/albedo.png} &
    \includegraphics[width=0.30\textwidth]{\assetroot/renders/view_002/gt_albedo.png}
    \\
    \textbf{2D Weiss} & \textbf{3D Weiss} & \textbf{Ground truth}
  \end{tabular}
  \caption{Gradient-space albedo estimation in image space and in the shared
  UV atlas, shown in view 2 and compared with ground truth.}
  \label{fig:weiss-albedo}
\end{figure}

In 2D, the estimate is sharper than the naive mean or median, but considerable
shading is still baked into the result. Extending the method to 3D improves the
separation further by combining surface-aligned evidence across views, yet
clear residual shading remains. A likely cause is the biased distribution of
the generated illumination; increasing both the number and diversity of
lighting conditions and views may reduce this bias.

\subsection{Locally Adaptive Temporal Constraints}

The Weiss estimator encodes a useful statistical prior: filter responses of
natural images are sparse, while illumination edges move between frames. This
implicitly favours piecewise-constant reflectance, since most reflectance
gradients are expected to be zero, but it does not explicitly decide where
local reflectance should be tied together. The much earlier Retinex family of
methods makes stronger use of local constancy, but does not exploit the extra
information available in an illumination sequence. We therefore evaluate
\href{https://openaccess.thecvf.com/content_cvpr_2015/papers/Laffont_Intrinsic_Decomposition_of_2015_CVPR_paper.pdf}{Laffont and Bazin, \emph{Intrinsic Decomposition of Image Sequences from Local
Temporal Variations} (2015)}, referred to as \texttt{laffont15} in our
implementation.

The method combines three types of constraints:

\begin{enumerate}
  \item \textbf{Adaptive local constraints.} Every frame is divided into
  overlapping $3\times3$ patches. If shading is constant within a patch, its
  pixel values differ from the patch reflectance only by a single scale. The
  method examines each patch across time and uses robust reweighting to reduce
  the influence of frames containing a shadow boundary, specularity, or other
  violation of locally constant shading.

  \item \textbf{Long-range pairwise constraints.} Local patches cannot connect
  regions separated by depth or normal discontinuities. Pixels are therefore
  clustered by their scale-normalised temporal appearance profiles, and
  distant pixels from the same cluster are paired. Proportional profiles
  indicate similar shading changes and constrain the relative reflectance of
  the two pixels. These constraints are also estimated robustly across time.

  \item \textbf{Scale regularisation.} Local and pairwise relations alone admit
  the trivial zero-reflectance solution and do not determine absolute scale. A
  very weak chromaticity term anchors the solution without dominating the
  temporal constraints.
\end{enumerate}

The three terms form one sparse least-squares problem, solved independently for
each colour channel:
\begin{equation}
  E(R) = E_{\mathrm{local}}(R)
       + \gamma_{\mathrm{pair}} E_{\mathrm{pair}}(R)
       + \gamma_{\mathrm{reg}} E_{\mathrm{reg}}(R).
  \label{eq:laffont15-energy}
\end{equation}

\begin{figure}[htbp]
  \centering
  \begin{tabular}{@{}c@{\hspace{1.2em}}c@{}}
    \includegraphics[width=0.40\textwidth]{\assetroot/albedo_2d/laffont15/view_002/albedo.png} &
    \includegraphics[width=0.40\textwidth]{\assetroot/renders/view_002/gt_albedo.png}
    \\
    \textbf{Laffont15} & \textbf{Ground truth}
  \end{tabular}
  \caption{Albedo estimated from the 15 relightings of view 2 using local
  temporal variations and long-range pairwise constraints.}
  \label{fig:laffont15-albedo}
\end{figure}

This approach removes the spatially varying shading remarkably well. Its main
failure in this example is instead the absolute reflectance scale: dark neutral
materials are reconstructed as much brighter than in the ground truth.
Normalising temporal appearance profiles removes their magnitude, so bright
white and near-black neutral patches can have almost identical profile
directions and be grouped together. Since both are also close to neutral
chromaticity, the weak scale regulariser cannot reliably separate them and can
pull the dark material towards brighter tones.

\subsubsection{Proposed Extension to 3D}

The 3D extension is still work in progress. One possible way to carry the local
part of the method into 3D would be to replace square image patches with, for
example, geodesic kernels on the mesh surface. Since geometry is available, we
could also inspect the surface normals and depth continuity and avoid applying
the local constraint across clearly non-planar or disconnected regions. The
constraints could then be estimated directly in shared surface coordinates
using observations from all views and illuminations.

There is also some flexibility in how the long-range pairs are chosen. Rather
than relying entirely on clusters of normalised appearance profiles, we could
ask a human or an image-understanding agent to suggest a small number of
high-confidence surface pairs that are likely to share a material or a shading
response. This may be a practical way to avoid incorrect links between dark
and bright neutral regions while still connecting parts of the surface that
are separated by geometric discontinuities.

\clearpage
\section{Next Steps}

\subsection{Other Albedo Estimation Methods}

I did not get around to trying all of the available methods. Two remaining
classical approaches and one learning-based approach could be particularly
relevant:

\begin{itemize}
  \item \href{https://people.csail.mit.edu/sparis/publi/2012/sigasia_intrinsic/Laffont_12_Coherent_Intrinsic_Images.pdf}{\emph{Coherent Intrinsic Images from Photo Collections}}
  is interesting because it is multi-view by design. It does not assume known
  geometry, but instead reconstructs it alongside the intrinsic decomposition
  and uses the resulting correspondences to obtain a coherent albedo across
  views.

  \item \href{https://www.microsoft.com/en-us/research/wp-content/uploads/2016/02/eccv2004_intrinsic.pdf}{\emph{Estimating Intrinsic Images from Image Sequences with Biased Illumination}}
  explicitly addresses biased lighting, such as sequences in which most light
  comes from above. This could be useful here because image-generation models
  are also likely to favour common, top-lit illumination. One caveat is that
  the paper does not report quantitative evaluation results, so its practical
  performance may vary in our setting.

  \item \href{https://www.cs.cornell.edu/projects/bigtime/}{\emph{Learning Intrinsic Image Decomposition from Watching the World}}
  provides a learning-based alternative. It learns from image sequences of the
  same scene under changing illumination and uses consistency across the
  sequence instead of requiring ground-truth intrinsic decompositions. A
  similar approach could make sense here because it would avoid direct albedo
  supervision, which could otherwise bias the model towards synthetic training
  data and defeat part of the motivation for this setup.
\end{itemize}

\subsection{More Consistent Relighting Sequences}

Another possible direction would be to fine-tune a video-to-video generation
model. We could condition it on a video of the baked-lighting RGB rendering as
the object rotates and ask it to produce a more consistent, light-stage-like
video. The current image-generation model produces each relighting one frame
at a time, so it is naturally more prone to small inconsistencies between
views. These inconsistencies are then difficult for the albedo-estimation
methods to resolve. I do not think this is the main limiting factor at the
moment, however, so this is probably a lower priority for now.

\end{document}
